\documentclass[10pt,a4paper]{amsart}
\usepackage[margin=19mm]{geometry}
\usepackage{amsmath,amssymb,mathtools}
\usepackage[scr=rsfs,cal=euler]{mathalfa}
\usepackage{xfrac}
\usepackage[unicode,hidelinks,bookmarks=false]{hyperref}
\newcommand{\ie}{i.e.\ }
\newcommand{\cf}{cf.\ }
\newcommand{\resp}{respectively\ }
\newcommand{\Subsection}{\subsection}
\providecommand{\nobreakdash}{\nobreak\hbox{-}}
\setlength{\emergencystretch}{2em}
\multlinegap0pt
\usepackage{amssymb}
\usepackage{mathtools}
\usepackage[shortlabels]{enumitem}
\usepackage{tikz}
\newtheorem{lem}{Lemma}
\newcommand{\Acal}{\mathcal A}
\newcommand{\R}{\mathbb R}
\newcommand{\dd}{\,d}
\newcommand{\loc}{\mathrm{loc}}
\title{Critical $p$-Laplace equations with monotone coefficients: Liouville classification and a Schoen-type Harnack inequality}
\author{Yi Ru-Ya Zhang}
\date{Source: arXiv:2608.09113v2 (CC BY 4.0)}
\begin{document}
\maketitle

\noindent\emph{Source and licence.} Critical $p$-Laplace equations with monotone coefficients: Liouville classification and a Schoen-type Harnack inequality. Version arXiv:2608.09113v2 (CC BY 4.0), distributed by arXiv under the Creative Commons Attribution 4.0 International licence, http://creativecommons.org/licenses/by/4.0/, as recorded in the arXiv licence metadata for this exact version and shown on the abstract page https://arxiv.org/abs/2608.09113v2. Full licence notice: \texttt{licenses/arxiv-2608.09113-CC-BY-4.0.txt}.\par\smallskip

\noindent\textit{Editorial scope.} Counted unit: the article's lem lem:unique-tangent (source0.tex lines 1854--1868). The article's complete original proof of it is delivered with it (source0.tex lines 1870--2045, one \texttt{proof} environment of 176 lines). No mathematics is written, repaired or re-derived: the statement and the proof are the article's own lines, in the article's own notation, and no retained line is substituted or re-flowed.  Retained uncounted as the source-local context the proof needs: lines 589--613; lines 690--692; lines 1838--1841.  Nothing was omitted from any retained span, so the package carries no omission record.  No retained line contains a citation, so the wrapper carries no bibliography and no bibliography entry.  No retained line contains a figure environment, an image-inclusion command or drawing code, so no image is delivered and the package declares no asset dependency.  Editorial formatting only, in addition: an AMS \texttt{amsart} preamble that restates the article's own declarations and macros used by the retained lines, namely the environments \texttt{lem} on separate counters, together with the standard packages those lines need.  The retained environments print in this document as Lemma 1 in order of appearance, whereas the article prints the same items with the numbering of its own full document.  The complete native source of the article as distributed in the arXiv e-print for this exact version is bundled unchanged as \texttt{sources/207215/source0.tex}.
\begin{lem}
\label{lem:uniform-C1-compactness}
Let $B_{2R}(x_0)\subset\mathbb{R}^n$, and let
$$
u_j \in W^{1,p}\bigl(B_{2R}(x_0)\bigr)\cap L^\infty\bigl(B_{2R}(x_0)\bigr)
$$
be weak solutions of the boundary value problem
$$
-\Delta_p u_j = f_j \qquad \text{in } B_{2R}(x_0),
$$
where the source terms $f_j$ are not assumed to have a fixed sign (i.e., they may take both positive and negative values). Assume that
$$
\sup_j\|u_j\|_{L^\infty(B_{2R}(x_0))}\le M_0,
\qquad \sup_j\|f_j\|_{L^\infty(B_{2R}(x_0))}\le F_0.
$$
Then there exist $\beta=\beta(n,p)\in(0,1)$ and
$C=C(n,p,R,M_0,F_0)$ such that
$$
\|u_j\|_{C^{1,\beta}(B_R(x_0))}\le C \qquad\text{for every }j.
$$
Consequently, $\{u_j\}$ is precompact in $C^1(B_R(x_0))$.
Moreover, if $u_j\to u$ in $C^1_{\rm loc}$ and
$f_j\to f$ in $L^1_{\rm loc}$, then $u$ is a weak solution of
$-\Delta_pu=f$.
\end{lem}

\begin{equation}\label{eq:transformed-equation}
wP_w=b_p|\nabla w|^p+c_\tau(w)
\end{equation}

\begin{equation}\label{eq:sharp-gradient-w}
b_p|\nabla w|^p \le w\int_1^w\frac{c_\tau(t)}{t^2}\dd t
\quad\text{in }\R^n.
\end{equation}

\begin{lem}\label{lem:unique-tangent}
Let $x_*\in\R^n$ satisfy $w(x_*)=1$, and define
\begin{equation}\label{eq:K-tau}
    K_\tau :=\frac1{p'}\left(\frac{c_{\tau,0}}n\right)^{1/(p-1)}.
\end{equation}
Then
\begin{equation}\label{eq:tangent-C1}
    \frac{w(x_*+rx)-1}{r^{p'}} \to K_\tau|x|^{p'} \qquad\text{in }C^1_{\loc}(\R^n)
\end{equation}
as $r\to 0$. In particular, there exists $r_0>0$ such that
$$
\nabla w(x)\ne0 \qquad \text{whenever }\ 0<|x-x_*|<r_0.
$$
Thus $x_*$ is the only critical point of $w$ in $B_{r_0}(x_*)$.
\end{lem}

\begin{proof}
Since $w\ge1$ and $w(x_*)=1$, the point $x_*$ is a
global minimum and $\nabla w(x_*)=0$.
Set $y(x):=w(x)-1.$
By continuity of $w$ and $c_\tau$, there exist $\rho,\delta,C>0$ such that
$$
0\le y\le\delta \quad\text{and}\quad
w\int_1^w\frac{c_\tau(t)}{t^2}\,\dd t \le Cy
\qquad\text{in }B_\rho(x_*).
$$
Hence \eqref{eq:sharp-gradient-w} gives
$$
|\nabla y|^p\le Cy \qquad\text{in }B_\rho(x_*).
$$

For $\epsilon>0$, define
$$
\Psi_\epsilon := (y+\epsilon)^{1/p'} -\epsilon^{1/p'}.
$$
Then
$$
|\nabla\Psi_\epsilon|^p
=\left(\frac1{p'}\right)^p (y+\epsilon)^{-1}|\nabla y|^p \le C\frac{y}{y+\epsilon} \le C.
$$
Since $\Psi_\epsilon(x_*)=0$, integration along line
segments in $B_\rho(x_*)$ gives
$$
\Psi_\epsilon(x)\le C|x-x_*|.
$$
Letting $\epsilon\to 0$, we obtain
\begin{equation}\label{eq:w-minimum-growth}
    0\le w(x)-1\le C|x-x_*|^{p'}
\end{equation}
for $x$ sufficiently close to $x_*$.

For $r>0$, define
$$
W_r(x):=\frac{w(x_*+rx)-1}{r^{p'}}.
$$
Then the transformed equation \eqref{eq:transformed-equation} becomes
\begin{equation}\label{eq:Wr-equation}
 \Delta_pW_r  =\frac{b_pr^{p'}|\nabla W_r|^p
+c_\tau(1+r^{p'}W_r)}{1+r^{p'}W_r}.
\end{equation}
The growth bound \eqref{eq:w-minimum-growth} bounds $W_r$ on each fixed ball.  Dividing  the scaled version of \eqref{eq:sharp-gradient-w} by $r^{p'}$, we obtain
\begin{equation}\label{eq:W-upper}
b_p  |\nabla W_r|^p \le \frac{1+r^{p'}W_r}{r^{p'}}\int_1^{1+r^{p'}W_r}\frac{c_\tau(t)}{t^2}\, \dd t \le C_R W_r     \qquad\text{on }\  B_R \ \text{ with }\  R>0
\end{equation}
for all sufficiently small $r$.
Hence, on every fixed ball $B_{2R}$, both $W_r$ and the right-hand side of \eqref{eq:Wr-equation} are uniformly bounded for all sufficiently small $r$.  Writing
$$
-\Delta_pW_r=f_r,
\qquad f_r:=-\frac{b_pr^{p'}|\nabla W_r|^p
+c_\tau(1+r^{p'}W_r)}{1+r^{p'}W_r},
$$
we have
$$
\sup_{0<r<r_R}\|f_r\|_{L^\infty(B_{2R})}<\infty,
$$
and Lemma~\ref{lem:uniform-C1-compactness} therefore makes
$\{W_r\}_{0<r<r_R}$ precompact in $C^1(B_R)$.

Take a sequence $r_j\to 0$ and let $W$ be a subsequential limit.  Then
\begin{equation}\label{eq:tangent-limit-system}
 \Delta_p W=c_{\tau,0},
 \qquad  W\ge0,  \qquad  W(0)=0,
 \qquad  b_p|\nabla W|^p\le c_{\tau,0}W  \quad \text{in } \ \R^n.
\end{equation}
Indeed, the equation follows from \eqref{eq:Wr-equation}. To obtain the last inequality, we divide \eqref{eq:sharp-gradient-w} by $r_j^{p'}$ and use the identity $c_\tau(1)=c_{\tau,0}$ to get that
$$
 \frac{1}{r_j^{p'}}\bigl(1 + r_j^{p'} W_{r_j}\bigr)
 \int_1^{1 + r_j^{p'} W_{r_j}} \frac{c_\tau(t)}{t^2}\,\dd t  \;\to\;  c_{\tau,0}\,W  \quad \text{locally uniformly.}
$$

The differential inequality in \eqref{eq:tangent-limit-system} needs to be used through an explicit zero-set regularization.  Towards this, we fix $\epsilon>0$ and define
$$
\Phi_\epsilon := (W+\epsilon)^{1/p'} -\epsilon^{1/p'}.
$$
Using the last inequality in
\eqref{eq:tangent-limit-system}, we obtain
\begin{align*}
|\nabla\Phi_\epsilon|^p &=\left(\frac 1 {p'}\right)^p (W+\epsilon)^{-1}|\nabla W|^p\\
&\le\left(\frac 1 {p'} \right)^p \frac{c_{\tau,0}}{b_p} \frac{W}{W+\epsilon} \le \left(\frac1{p'}\right)^p \frac{c_{\tau,0}}{b_p}.
\end{align*}
Since $\Phi_\epsilon(0)=0$, the Lipschitz estimate stated above implies that
$$
\Phi_\epsilon(x) \le \frac{1}{p'}\left(\frac{c_{\tau,0}}{b_p}\right)^{1/p}\lvert x\rvert.
$$
Equivalently, by raising both sides of this inequality to the power $p'$, we obtain
$$
\lvert \Phi_\epsilon(x)\rvert^{p'} \le \left(\frac{1}{p'}\right)^{p'}\left(\frac{c_{\tau,0}}{b_p}\right)^{1/(p-1)}\lvert x\rvert^{p'}.
$$
Letting $\epsilon\to 0$ gives
\begin{equation}\label{eq:tangent-upper-cone}
W(x)\le K_\tau|x|^{p'}.
\end{equation}
The constant in \eqref{eq:tangent-upper-cone} is exactly
\eqref{eq:K-tau} as $b_p=n/p'$.

We next show that the upper bound is indeed attained. To begin with, note that the  gradient estimate in \eqref{eq:tangent-limit-system} and \eqref{eq:tangent-upper-cone} imply
\begin{equation}\label{eq:tangent-stress-bound}
|\Acal(\nabla W(x))|=|\nabla W|^{p-1}\le \left(\frac{p'}{n} c_{\tau,0}\right)^{\frac {p-1}{p}} W^{\frac {p-1}{p}} \le\frac{c_{\tau,0}}n|x| \qquad\text{for almost every } x\in \R^n.
\end{equation}

For every $R>0$, let
$$
\nu(x):=\frac{x}{R} \qquad (x\in\partial B_R)
$$
denote the outward unit normal.
Testing the weak equation \eqref{eq:tangent-limit-system} with radial cutoffs converging to $\mathbf1_{B_R}$ and using the coarea formula yields that, for almost every $R>0$,
\begin{equation}\label{eq:tangent-flux}
 \int_{\partial B_R}\Acal(\nabla W)\cdot\nu\dd \mathcal H^{n-1}
 =c_{\tau,0}|B_R|  =\frac{c_{\tau,0}}nR|\partial B_R|.
\end{equation}


Observe that, on $\partial B_R$, \eqref{eq:tangent-stress-bound} gives
$$
\Acal(\nabla W)\cdot\nu \le |\Acal(\nabla W)|\le \frac{c_{\tau,0}}nR.
$$
Since the integral in \eqref{eq:tangent-flux} equals the
integral of the last upper bound, the nonnegative function
$$
\frac{c_{\tau,0}}nR -\Acal(\nabla W)\cdot\nu
$$
has zero integral over $\partial B_R$.  Hence
$$
\Acal(\nabla W)\cdot\nu =\frac{c_{\tau,0}}nR
$$
for $\mathcal H^{n-1}$-almost every point of
$\partial B_R$.
Therefore, the equality must hold in both inequalities above. The equality case of the Cauchy--Schwarz inequality gives
$$
\Acal(\nabla W) = \frac{c_{\tau,0}}nR\,\nu
= \frac{c_{\tau,0}}n x
$$
almost everywhere on $\partial B_R$.
Since this holds for almost every $R>0$, polar-coordinate
Fubini gives
\begin{equation}\label{eq:tangent-stress-identity}
\Acal(\nabla W) =\frac{c_{\tau,0}}n x \qquad \text{ almost everywhere in }\R^n.
\end{equation}

 
Moreover, thanks to 
$$
\Acal(\xi)=|\xi|^{p-2}\xi, \qquad \Acal^{-1}(\zeta)=|\zeta|^{p'-2}\zeta,
$$
the identity \eqref{eq:tangent-stress-identity} yields
$$
\nabla W(x) =\left(\frac{c_{\tau,0}}n\right)^{1/(p-1)}
    |x|^{p'-2}x
$$
almost everywhere.  As $W\in C^1(\R^n)$, this identity actually holds everywhere.  Integrating along the segment from $0$ to $x$, and using $W(0)=0$, we finally arrive at
$$
W(x) = \frac1 {p'} \left(\frac{c_{\tau,0}}n\right)^{1/(p-1)} |x|^{p'} = K_\tau|x|^{p'}.
$$
Since every convergent subsequence has the same limit and the
family $\{W_r\}$ is locally precompact in $C^1$, the
convergence in \eqref{eq:tangent-C1} holds for the full family
as $r\to 0$.

Finally, on the annulus
$$
    A:=\{x\in\R^n:1/2\le|x|\le2\},
$$
the gradient of $K_\tau|x|^{p'}$ is bounded away from zero.
Hence the $C^1(A)$-convergence gives
$\nabla W_r\ne 0$ on $A$
for every sufficiently small $r$.
Then rescaling back to $w$ shows that
$$
 \nabla w(x)\ne 0 \qquad \text{whenever }0<|x-x_*|<r_0
$$
for some $r_0>0$.
\end{proof}

\end{document}
